\chapter{Bevezetés}

\section{Motiváció és a témaválasztás indoklása}

\section{A megoldandó feladat rövid, közérthető leírása}

\section{A dolgozat célja és hozzájárulásai}

\section{A dolgozat felépítése}

\chapter{Felhasználói dokumentáció}

\section{A megoldott probléma rövid megfogalmazása}

\section{A felhasznált módszerek rövid áttekintése}

\section{Rendszerkövetelmények}

\section{Telepítés és első indítás}

\section{A grafikus felhasználói felület bemutatása}

\section{Tipikus felhasználási forgatókönyvek}

\section{Hibaüzenetek és gyakori problémák}

\chapter{Fejlesztői dokumentáció}
Ebben a fejezetben a megvalósítandó szoftver célját és funkcióit formális, matematikai keretben definiálom, mivel a későbbi fejezetekben tárgyalt 
algoritmusok (FGSM, PGD) és a védekezési modul is ugyanerre a jelölésrendszerre épülnek majd. 
Először a probléma szereplőit, a modellt, a bemenetet és a hozzá tartozó címkét vezetem be, majd ezek segítségével pontosan megfogalmazom, mit 
is jelent az ellenfél alapú perturbáció. 
Ezután az itt bevezetett jelölésrendszerre épülve sorra bemutatom a két megvalósított támadási algoritmust (FGSM, PGD), a mögöttük álló tervezési 
döntéseket, a kapcsolódó munkákat, a védekezési stratégiákat, végül a szoftver architektúráját, a felhasznált technológiákat és az egyes 
komponensek megvalósítását. 

\section{A probléma részletes specifikációja}

\subsection{A modell, a bemenet és a címke}
Legyen adott egy $f_\theta$ neurális háló alapú osztályozó modell, melynek $x$ egy bemenete és $y$ annak helyes címkéje. 
Az $x$ bemenet ebben az esetben egy kép pixeleinek valamely három színű kódolása, vektoriális formában,
az $y$ pedig a kép helyes osztályának indexe, például $y=42$ az agámák osztályát reprezentálja. 
Amennyiben a képhez tartozó $y$ helyes címke nem ismert, úgy a modell által meghatározott legvalószínűbb címkét $\hat{y}$-nal jelöljük, és ezt 
használjuk referenciacímkeként. 

\subsection{A támadás célja}
Célunk az $x$ képre egy $x_{\mathrm{adv}}$ bemenet előállítása, mely az eredeti $x$-től szemmel csak alig látható perturbációban különbözik,
viszont a modell által adott $\hat{y}_{\mathrm{adv}} = \arg\max_i f_\theta(x_{\mathrm{adv}})_i$ predikció a referenciacímkétől eltér, vagy a 
klasszifikáció elbizonytalanodik. 

\subsection{Targeted és non-targeted mód}
Alapvetően két támadási módot különböztetünk meg, az ún. targeted és non-targeted módot. 
Non-targeted esetben a támadás célja minél jobban elbizonytalanítani a helyes osztályozást, tehát hogy agáma helyett bármi mást mutasson. 
Targeted módban azt szeretnénk, hogy egy kijelölt $y_{\mathrm{target}}$ osztály irányába változzon a képünk (pl.\ a komodói varánusz osztálya felé). 

\subsection{A white-box hozzáférési feltétel}
A támadás feltétele, hogy az $f_\theta$ modell gradienseihez és paramétereihez teljes hozzáférésünk legyen. 
Ezt white-box hozzáférésnek nevezzük az AI red teaming kontextusában. 

\subsection{A szoftver célja és beállítható paraméterei}
E szoftver célja a fenti ellenfél alapú támadó tesztkörnyezet szoftveres megvalósítása, különböző támadási algoritmusokkal (FGSM, PGD). 
A szoftver a PyTorch (torchvision) előtanított modelljeire épül, mely architektúra tetszőleges torchvision-modellel
bővíthető. Jelenleg kettő van integrálva (ResNet50, MobileNetV2). Beállíthatóak a futás paraméterei:
\begin{itemize}
  \item modell: ResNet50 vagy MobileNetV2, amelyen a támadást végrehajtjuk,
  \item $\epsilon$, mely meghatározza a perturbáció maximális mértékét, tehát azt, hogy egy pixelen legfeljebb mekkora
        változást tapasztalhatunk, azaz a változás felső határát,
  \item célosztály, $y_{\mathrm{target}}$, ha a mód targeted,
  \item iterációszám (PGD esetén hány lépésig engedjük futtatni az algoritmust).
\end{itemize}

\subsection{A felhasználói funkciók vázlatos áttekintése}

A fenti specifikációt a szoftver az alábbi, a felhasználó számára közvetlenül elérhető funkciókkal valósítja meg:
\begin{itemize}
  \item a fájlrendszerről kiválasztott képet legördülő menüből kiválasztott modellel klasszifikáljuk,
  \item támadási algoritmus és paraméterek meghatározásával perturbációt és ellenfél alapú képet generálunk,
  \item megjelenítjük az eredeti képet, a perturbációt és az ellenfél alapú képet, továbbá mindkét (eredeti és
        ellenfél alapú) képet osztályozzuk is,
  \item a védekezési modulban JPEG-tömörítést vagy más képfeldolgozási szűrőt alkalmazhatunk az ellenfél alapú perturbáció 
  hatásának utólagos csökkentésére vagy semlegesítésére.
\end{itemize}

\section{Elméleti háttér és a használt fogalmak definíciója}

\subsection{Az osztályozó modell és a logitok}

Jelöljük $f_\theta(x) = z$-vel a modell outputjaként előállt valós számok tömbjét, ahol $z \in \mathbb{R}^C$ és $C$ az osztályok száma (ImageNet esetén $C=1000$). 
Ezt logitvektornak is nevezzük, mely elnevezés a logisztikus regresszióból származik, ahol a $\mathrm{logit}(p) = \ln\frac{p}{1-p}$ függvény egy $p \in (0, 1)$ valószínűséget 
képez le a valós számegyenesre. 
A $z_i$ komponensek páronkénti különbsége éppen az $i$-edik és $j$-edik osztály esélyhányadosának logaritmusával egyezik meg, 
$z_i - z_j = \ln\frac{p_i}{p_j}$, innen származik az elnevezés.

\subsection{Valószínűségek és a modell predikciója}

Mivel a $z$ logitok önmagukban nem valószínűségek, a $\mathrm{softmax}$ függvénnyel alakítjuk át őket valószínűségi eloszlássá: 
\[
  p_i = \mathrm{softmax}(z)_i = \frac{e^{z_i}}{\sum_{j=1}^{C} e^{z_j}}, \qquad i = 1, \ldots, C.
\]
Tehát $p_i$ a modell szerint annak a valószínűsége, hogy $x$ az $i$-edik osztályba tartozik. 
A modell által meghatározott legvalószínűbb osztály indexét a legnagyobb valószínűségű komponens alapján kapjuk: 
\[
  \hat{y} = \arg\max_{i} p_i.
\]

\subsection{Veszteségfüggvény és bemeneti gradiens}

A $J(\theta, x, y)$ veszteségfüggvény a modell predikciója és a megadott helyes címke közötti eltérést méri. 
Esetünkben ez a keresztentrópia-veszteség 
\[
J(\theta, x, y) = -\log \frac{e^{z_y}}{\sum_{i=1}^{C} e^{z_i}},
\]
ahol $z = f_\theta(x)$ a logitvektor és $y$ a helyes osztály indexe. 
A veszteségfüggvény súlyok szerinti gradiense, $\nabla_\theta J(\theta, x, y)$, az a gradiensvektor, amelyet a backpropagation során a tanítás közben használunk a súlyok módosítására. 
Itt a gradiensszámítást egy további lépéssel a bemeneti rétegig visszük -- azaz a backpropagation folyamatát a hálón teljesen végigvezetjük --, és megkapjuk a bemeneti vektorra vonatkozó érzékenységet, $\nabla_x J(\theta, x, y)$. 
Ez a mennyiség azt mutatja meg, hogy egy adott bemeneti pixel vagy jellemző változása milyen mértékben befolyásolja a veszteséget, ezért ennek a mennyiségnek központi szerepe 
lesz a későbbiekben bemutatott, ellenfél alapú példát generáló algoritmusok leírásában. 

\subsection{Perturbáció és normakorlát}

Az ellenfél alapú képek definíciójánál szokás a perturbáció nagyságát valamely $\ell_p$ normával jellemezni; ebben a dolgozatban a $\ell_\infty$ normát használjuk. 
Egy perturbált kép általában az
\[
  x_{\mathrm{adv}} = x + \delta
\]
formában írható le, ahol $\delta$ a zaj, vagyis a képpontokban okozott változás. 
A $\ell_\infty$ normát a 
\[
  \|\delta\|_\infty = \max_i |\delta_i|
\]
kifejezés definiálja, azaz a legnagyobb abszolút értékű komponens adja a norma értékét. 
A perturbáció nagyságát ezzel a normával korlátozzuk: 
\[
  \|\delta\|_\infty \leq \epsilon.
\]
Ez esetünkben nagyon kézenfekvő, mivel a bemeneti kép minden pixelén (minden $\delta_i$ komponensen) külön-külön korlátozza a megengedett változást. 

\section{Ellenfél alapú támadások áttekintése}
A mély neurális hálók, köztük a képfelismerő modellek is, különösen érzékenyek a bemeneti adatok kicsi, ám célzott módosításaira~\cite{goodfellow2015explaining}. 
Az ilyen módosításokat ellenfél alapú perturbációknak nevezzük. 
Ezeknek a támadásoknak két algoritmusát fogjuk itt átvenni, melyek a taxonómiában a korábban említett white-box támadások közé tartoznak. 
A következő fejezetekben egy egylépéses ellenfél alapú támadást mutatunk be (FGSM), utána erre építve egy többlépéses algoritmust (PGD). 

\section{FGSM (Fast Gradient Sign Method)}
Az FGSM algoritmust Goodfellow és szerzőtársai a Google-nél végzett kutatásaik alapján javasolták~\cite{goodfellow2015explaining}. 
Az előző fejezetben megjegyeztem, hogy a $\nabla_x J(\theta, x, y)$ veszteségfüggvény gradiense az input értékre való érzékenységet jelenti. 
Természetes ötlet, hogy a bemenetet a gradienssel ellentétes irányba mozdítsuk el, ami várhatóan növeli a veszteségfüggvény értékét, azaz rontja a 
helyes osztályra adott predikciót. 
Ehhez feltesszük, hogy a modell az $x$ pont kellően kis környezetében közelítőleg lineárisan viselkedik -- ez Goodfellow és társai linearitási 
elmélete, amely szerint a modern neurális hálók, részben a szakaszonként lineáris aktivációk (pl.\ ReLU) miatt, lokálisan jól közelíthetők egy 
lineáris függvénnyel, még ha globálisan erősen nemlineárisak is~\cite{goodfellow2015explaining}. 
Az FGSM algoritmusban az elmozdítást egyszerűen csak az előjel-függvény ($\mathrm{sign}$) $\epsilon$-szorosával tesszük meg, tehát nem vesszük 
figyelembe a gradiens nagysága által adott érzékenységet, csak az irányát. 
Ily módon az $\epsilon \cdot \mathrm{sign}(\cdot)$ minden egyes $x_i$ komponensre kihasználja a teljes $\epsilon$ költségvetésünket. 
Tehát az FGSM ellenfél alapú példát a következő módon kapjuk meg: 
\[
  x_{\mathrm{adv}} = x + \epsilon \cdot \mathrm{sign}\big(\nabla_x J(\theta, x, y)\big).
\]

\subsection{Célzott változat}
Az alapváltozat a megadott helyes osztály veszteségét igyekszik maximalizálni, tehát minél inkább mást mutasson, mint ami a képen van. 
A célzott változatnál egy $y_{\mathrm{target}}$ célosztályt adunk meg, és azt szeretnénk, hogy a modell ebbe az osztályba sorolja az ellenfél alapú képet. 
\[
  x_{\mathrm{adv}} = x - \epsilon \cdot \mathrm{sign}\big(\nabla_x J(\theta, x, y_{\mathrm{target}})\big).
\]
Tehát a képlet csak annyiban változik, hogy távolodás helyett az $y_{\mathrm{target}}$-hoz tartozó veszteség minimuma felé közeledni szeretnénk, és $y$ helyett $y_{\mathrm{target}}$-ot használjuk. 

\subsection{Előnyök és korlátok}
Mivel a gradienst a tanításban is használt backpropagation segítségével hatékonyan tudjuk kiszámolni, ezért a támadás gyors. 
Hátránya, hogy csak egyetlen, durva lépést tesz a megfelelő irányba, ezért -- főleg az iteratív módszerekhez, például a PGD-hez képest -- 
kevésbé hatékonyan képes megtéveszteni az osztályozót. 

\section{PGD (Projected Gradient Descent)}
Ezt a módszert Madry és társai vezették be~\cite{madry2018towards}, és lényegében az FGSM egylépéses ötletének iteratív általánosítása. 
Nem egyetlen nagy lépést teszünk a gradiens előjele irányába, hanem ugyanabba az irányba több, kisebb lépést hajtunk végre egymás után. 
A motiváció emögött az, hogy az FGSM linearitási feltevése csak $x$ kellően kis környezetében állja meg a helyét, ezért előnyösebb, ha $T$ darab 
kisebb, $\alpha$ nagyságú lépéssel dolgozunk, és a gradienst minden lépésben újraszámoljuk. 
Legyen $S = \{\delta : \|\delta\|_\infty \leq \epsilon\}$ a megengedett perturbációk halmaza; ekkor $x+S$ adja meg az $x$ körüli, $\epsilon$ 
sugarú gömbbe eső, engedélyezett bemenetek halmazát. 
($\ell_\infty$ norma esetén ez geometriailag valójában egy kocka, de általános $\ell_p$ értelemben vett gömbként hivatkozunk rá.) 
Mivel az iteráció túlléphet $x+S$ korlátján, az algoritmus minden lépés után egy projekcióval visszavetíti az eredményt erre a halmazra. 
Ez FGSM esetén a $\mathrm{sign}$ függvény miatt automatikusan teljesült, itt viszont $T$ és $\alpha$ megválasztásától függően külön biztosítani kell. 
Az algoritmus formális definíciója tehát 
\[
  x^{t+1} = \Pi_{x+S}\big(x^t + \alpha \cdot \mathrm{sign}(\nabla_x J(\theta, x^t, y))\big), \qquad x^0 = x,
\]
ahol $\Pi_{x+S}$ az $x+S$ halmazra való projekció.

\subsection{Célzott változat}
Az FGSM-hez hasonlóan itt is definiálunk egy célzott változatot, amely az ottanival analóg módon adható meg: 
\[
  x^{t+1} = \Pi_{x+S}\big(x^t - \alpha \cdot \mathrm{sign}(\nabla_x J(\theta, x^t, y_{\mathrm{target}}))\big).
\]

\subsection{Véletlen kiindulási pont}
Madry és társai kísérletileg mutatták meg, hogy a veszteségfelszín lokális maximumai jól koncentráltak, így egy véletlen kezdőpontból indított PGD 
is megbízhatóan jó megoldást talál~\cite{madry2018towards}. 
Ennek a modellek robusztusságának tesztelésében is van szerepe: mivel a veszteségfelszínnek több, egymástól távol eső lokális maximuma is lehet, egy 
fix (pl.\ magából $x$-ből induló) kezdőpont esetén a szerencsének nagy szerepe lehetne, és az eredmény nem feltétlenül mutatná a modell valódi 
sebezhetőségét. 
Ezért a véletlen kezdőpontot az $\epsilon$ sugarú gömbön egyenletes eloszlással választjuk meg: 
\[
  x^0 = \Pi_{x+S}(x + \eta), \qquad \eta \sim \mathcal{U}(-\epsilon, \epsilon)^d.
\]

\subsection{Az $\alpha$ és $T$ paraméterek megválasztása}
A PGD viselkedését az $\epsilon$ perturbációs korláton felül két további paraméter határozza meg. 
$T$ az iterációk számát jelöli, ez adja meg az algoritmus számítási költségét, ugyanakkor több iteráció jellemzően erősebb támadást eredményez. 
Az $\alpha$ lépésnagyságot úgy kell megválasztani, hogy túl kicsi érték esetén lassú a konvergencia, túl nagy érték esetén pedig instabilitáshoz 
vezethet. 
Tipikus választás: $\alpha = \epsilon / T$ vagy $\alpha = 2{,}5 \cdot \epsilon / T$.

\subsection{Előnyök és korlátok}
A PGD az FGSM-nél lényegesen erősebb támadásokat képes generálni, mivel több kis lépésben közelíti a lokális optimumot. 
A véletlen kezdőponttal kombinálva különösen hatékony a robusztusság kiértékelésében. 
Hátránya, hogy $T$-szer annyi gradiens-számítást igényel, mint az FGSM, így nagyobb a számítási költsége. 

\section{Továbbmutató elméleti vonatkozások}
Ebben a fejezetben azokat az elméleti vontkozásokat foglalom össze, amelyek nem képezik szigorúan a dolgozatban megvaósított szoftver részét, ám fontos 
kiegészítői annak.

\subsection{Transzferálhatóság}
% TODO: az ellenfél alapú példák azon tulajdonsága, hogy egy modellre generált perturbáció gyakran más,
% eltérő architektúrájú vagy más adathalmazon tanított modelleket is megtéveszt -- ez a jelenség
% teszi lehetővé a black-box támadásokat, és mindkét cikk (Goodfellow, Madry) tárgyalja.

\subsection{Ellenfél alapú tanítás}
Ezeknek a támadási algoritmusoknak a szerepe nem csak a modell megtévesztése, hanem felhasználhatók a tanítás során a modell robusztusabbá tételére is. 
Ennek lényege, hogy a tiszta $x$ kép mellett a hozzá tartozó $x_{\mathrm{adv}}$ ellenfél alapú képet is bevonjuk a veszteségfüggvénybe, és a modell 
paramétereit az így kapott kombinált célfüggvényre optimalizáljuk. 

Goodfellow és társai az FGSM-mel előállított $x_{\mathrm{adv}}$ perturbációra a következő kombinált veszteségfüggvényt 
javasolták~\cite{goodfellow2015explaining}: 
\[
\tilde{J}(\theta, x, y) = \lambda \, J(\theta, x, y) + (1-\lambda)\, J(\theta, x_{\mathrm{adv}}, y), \qquad \lambda = 0{,}5.
\]
Az új veszteségfüggvény tehát a tiszta és az ellenfél alapú minta hibáját egyaránt bünteti, méghozzá úgy, hogy $x_{\mathrm{adv}}$-ot minden tanítási lépésben a 
modell aktuális $\theta$ paramétereivel újraszámolják. 
Így tehát a modell a tanítás során folyamatosan fejlődik, és az aktuális állapota által generált ellenfél alapú példákon tanul tovább dinamikusan, nem pedig 
egy előre legyártott, statikus ellenfél alapú halmazon. 
Kísérleteik megmutatták, hogy ez egy hatékony módszer: az FGSM-mel előállított ellenfél alapú példákon mért hibaarány 89,4\%-ról 17,9\%-ra csökkent. 

Madry és társai megmutatták, hogy PGD-vel is alkalmazható az ellenfél alapú tanítás, és erősebb védekezést eredményez~\cite{madry2018towards}. 
Itt azonban a tiszta és ellenfél alapú veszteséget kombináló FGSM-formulával ellentétben nincs szó keverésről: a tanítás során a tiszta $x$ kép 
helyett kizárólag a hozzá tartozó, PGD-vel előállított $x_{\mathrm{adv}}$-on számolják a veszteséget, tehát a modell csak az ellenfél alapú bemeneteken tanul. 
Mivel a PGD erősebb támadás, mint az FGSM, az ellene edzett modell a gyengébb, egylépéses támadásokkal szemben is ellenállóbb marad, míg a fordítottja nem 
feltétlenül igaz. 

\subsection{Linearitás}
% TODO: a linearitási hipotézis részletesebb kifejtése -- miért viselkednek a modern neurális hálók
% (ReLU, maxout stb. miatt) lokálisan lineárisan, és ez hogyan magyarázza az ellenfél alapú példák
% létezését magas dimenziószámú bemeneteken (Goodfellow et al. levezetése lineáris modellekre).

%\subsection{Saddle point}
% TODO: Madry et al. min-max (saddle point) optimalizálási kerete: a belső maximalizálás a
% legrosszabb esetű perturbáció megtalálása (támadás), a külső minimalizálás a robusztus modell
% tanítása (védekezés) -- ez adja a PGD elméleti indoklását mint "univerzális első rendű ellenfél".

\subsection{Kapcsolódó munkák és további támadási módszerek}

\section{Tervezési döntések és indoklásuk}

\section{Védekezési stratégiák áttekintése}

\section{A rendszer architektúrája}

\section{A felhasznált technológiák}

\section{Az egyes komponensek megvalósítási részletei}

\chapter{Tesztelés}

\section{Egységtesztek}

A program tesztkészletének felépítésénél a szakirodalomból ismert \textit{tesztpiramisnak} nevezett
hierarchiát követtem~\cite{fowler2012testpyramid}. Ennek alapját a gyorsabb és kevés komponenst érintő tesztek, az egységtesztek képezik, 
melyek a lehető legkisebb elkülöníthető egységek tesztelését valósítják meg. Futási idejük 
jellemzően ezredmásodperces tartományban mozog és akár lefedhetik az input paraméterek teljes értékkészletét. Az elkészítésüknél fontos szempont a 
többi egységtől való függetlenség (izoláció) és az perifériáktól való függetlenség, I/O, adatbázis, hálózat, stb. Ezeknek a teszteknek az elkészítésére 
érdemes nagy figyelmet fordítani, mivel a korai szakaszban feltárt hibák javítása általában lényegesen olcsóbb, mint a később, integráció vagy éles 
környezetben felfedezetteké.

Az egységtesztek tudatosan nem a valódi, előtanított ResNet50/MobileNetV2 modelleket használják. Az
\texttt{attack\_engine} teszteknél egy \texttt{tiny\_model} fixture ad egy $5$ osztályos,
$8\times8$ pixeles bemenetet feldolgozó, egyetlen lineáris rétegből álló hálót. Ez a döntés a tesztek sebességét 
(a valódi háló letöltése és felépítése másodperces nagyságrendű lenne) és izoláltságát szolgálja: a teszt kizárólag a 
saját kódunk logikáját vizsgálja, nem a harmadik féltől származó architektúra pontosságát. 

Minden véletlent felhasználó fixture (pl. a bemeneti képet generáló \texttt{sample\_image\_tensor})
maga hívja meg a \texttt{torch.manual\_seed}-et közvetlenül az adat előállítása előtt, így a
determinizmus független a fixture-ök végrehajtási sorrendjétől. Az egységtesztek szintjén nem támadási sikerrátát mérünk (ehhez sok mintás, statisztikai kiértékelés
kellene), hanem az algoritmusok tulajdonságait ellenőrizzük
konkrét, kis méretű bemeneteken.

Ebben a részben a számosságuk miatt nem mutatom be az összes egységtesztet, csak a középpontban álló \texttt{attack\_engine.py} 
modult mivel ez valósítja meg az alkalmazás lényegi logikáját (az FGSM és PGD algoritmusokat,
valamint az ezeket kiszolgáló \texttt{predict()} és \texttt{attack()} függvényeket). 

\bigskip
\noindent\textbf{A \texttt{predict()} és a diszpécser (\texttt{attack()}) tesztjei}

\medskip
\noindent\texttt{test\_predict\_re\-turns\_topk\_shapes}: Ellenőrzi, hogy \texttt{topk=3} esetén a visszaadott
valószínűség- és index-tenzor alakja $(1, 3)$, azaz a batch- és $k$-dimenzió megfelel a kérésnek.

\medskip
\noindent\rule{\linewidth}{0.3pt}

\medskip
\noindent\texttt{test\_predict\_re\-turns\_valid\_distribution}: A \texttt{topk=5} valószínűségek mindegyike a
$[0, 1]$ intervallumba esik, és -- mivel az 5 osztályos \texttt{tiny\_model} minden osztálya lekérdezésre
kerül -- összegük $1{,}0$-ra \texttt{allclose}. Ez a softmax kimenet valódi valószínűségi eloszlás voltát
igazolja.

\medskip
\noindent\rule{\linewidth}{0.3pt}

\medskip
\noindent\texttt{test\_predict\_top\-k\_returns\_highest\_first}: A visszaadott valószínűségek csökkenő
sorrendben szerepelnek ($\mathrm{probs}[0, i] \geq \mathrm{probs}[0, i+1]$), vagyis a \texttt{topk} valóban a
legnagyobb konfidenciájú osztályokat adja vissza, méghozzá rendezetten.

\medskip
\noindent\rule{\linewidth}{0.3pt}

\medskip
\noindent\texttt{test\_attack\_unk\-nown\_method\_raises}: Határeset: egy nem létező módszernév
(\texttt{"NEM VALID"}) megadásakor \texttt{ValueError}-t vár, mielőtt bármilyen gradiensszámítás történne --
a bemenet-validáció a drága számítás előtt fut le.

\medskip
\noindent\rule{\linewidth}{0.3pt}

\medskip
\noindent\texttt{test\_attack\_fgsm\_ig\-nores\_iterations\_argument}: Annak ellenére, hogy a hívó
\texttt{iterations=99}-et ad át, FGSM esetén a visszaadott \texttt{AttackResult.iterations\_run} értéke $1$ --
igazolva, hogy az FGSM valóban egylépéses, és a paramétert a diszpécser figyelmen kívül hagyja erre a
módszerre.

\medskip
\noindent\rule{\linewidth}{0.3pt}

\medskip
\noindent\texttt{test\_attack\_pgd\_re\-ports\_requested\_iteration\_count}: PGD esetén a kért iterációszám
($4$) változatlanul megjelenik az \texttt{iterations\_run} mezőben, alátámasztva, hogy a diszpécser a
paramétert továbbadja a \texttt{\_pgd} implementációnak.

\medskip
\noindent\rule{\linewidth}{0.3pt}

\medskip
\noindent\texttt{test\_attack\_fgsm\_dis\-patches\_only\_to\_fgsm}: \texttt{mocker.patch(..., wraps=\_fgsm)} és
\texttt{wraps=\_pgd} spy-okkal figyeli a hívásokat: \texttt{method="FGSM"} esetén
\texttt{\_fgsm.assert\_called\_once()} és \texttt{\_pgd.assert\_not\_called()} teljesül. A \texttt{wraps}
paraméter miatt a valódi implementáció is lefut, tehát az \texttt{AttackResult} összeállítása is ellenőrzött
marad.

\medskip
\noindent\rule{\linewidth}{0.3pt}

\medskip
\noindent\texttt{test\_attack\_pgd\_dis\-patches\_only\_to\_pgd}: Az előző teszt tükörképe:
\texttt{method="PGD"} esetén kizárólag \texttt{\_pgd} hívódik meg, \texttt{\_fgsm} nem. A két teszt együtt
zárja ki, hogy a diszpécser valamelyik ágon ,,átcsorogva'' mindkét algoritmust lefuttassa.

\bigskip
\noindent\textbf{Az FGSM implementáció (\texttt{\_fgsm}) tesztjei}

\medskip
\noindent\texttt{test\_fgsm\_targeted\_steps\_opposite\_direction\_of\_untargeted}: Azonos bemeneten és
$\epsilon$ mellett a targeted és non-targeted perturbáció pontosan ellentétes előjelű
($x_{\text{untargeted}} - x = -(x_{\text{targeted}} - x)$), ami közvetlenül a matematikai definícióból (a
gradiens előjelének megfordításából) következik.

\medskip
\noindent\rule{\linewidth}{0.3pt}

\medskip
\noindent\texttt{test\_fgsm\_clamps\_to\_valid\_pixel\_range}: Szélsőséges bemenet (csupa nulla pixel) és
nagy $\epsilon=0{,}5$ mellett is a kimenet minden eleme a $[0, 1]$ tartományban marad -- igazolva, hogy a
\texttt{clamp} lépés ténylegesen érvényesül, nem csak elméletben szerepel a specifikációban.

\medskip
\noindent\rule{\linewidth}{0.3pt}

\medskip
\noindent\texttt{test\_fgsm\_re\-sult\_is\_de\-tached\_from\_autograd}: A visszaadott tenzor
\texttt{requires\_grad} attribútuma \texttt{False}.

\medskip
\noindent\rule{\linewidth}{0.3pt}

\medskip
\noindent\texttt{test\_fgsm\_re\-sult\_has\_no\_graph\_history}: Szigorúbb bizonyíték az előzőnél:
\texttt{grad\_fn is None} \emph{és} \texttt{is\_leaf is True} egyszerre teljesül, ami kizárja azt az esetet
is, amikor egy tenzor \texttt{requires\_grad} értéke hamis, de mégis egy számítási gráf levele marad rejtett
hivatkozásokkal.

\medskip
\noindent\rule{\linewidth}{0.3pt}

\medskip
\noindent\texttt{test\_fgsm\_ze\-ro\_epsi\-lon\_re\-turns\_original}: Határeset: $\epsilon=0$ esetén a perturbáció
nulla, a kimenet \texttt{allclose} az eredeti képpel -- az algoritmus degenerált esetben identitásként
viselkedik.

\medskip
\noindent\rule{\linewidth}{0.3pt}

\medskip
\noindent\texttt{test\_fgsm\_un\-targeted\_re\-duces\_confidence\_in\_true\_label}: Viselkedési bizonyíték: a
\texttt{predict()}-en keresztül lekérdezett, a valódi címkéhez tartozó konfidencia a támadás \emph{után}
szigorúan kisebb, mint előtte -- ez nem csak a perturbáció létezését, hanem annak tényleges hatását igazolja
a modell kimenetére.

\medskip
\noindent\rule{\linewidth}{0.3pt}

\medskip
\noindent\texttt{test\_fgsm\_tar\-geted\_in\-creases\_confidence\_in\_target\_label}: A targeted változat
tükörpárja: a mesterségesen kiválasztott célosztály konfidenciája nő a támadás hatására, alátámasztva, hogy a
targeted mód valóban a célosztály felé tereli a predikciót, nem csupán ,,valamerre'' téveszti meg a modellt.

\bigskip
\noindent\textbf{A PGD implementáció (\texttt{\_pgd}) tesztjei}

\medskip
\noindent\texttt{test\_pgd\_re\-sult\_stays\_within\_epsilon\_ball}: Mérsékelt lépésméret ($\alpha=0{,}02$)
mellett a végső perturbáció minden pixelre $|\delta| \leq \epsilon + 10^{-6}$ (a tolerancia a lebegőpontos
kerekítést kompenzálja) -- ez az $\ell_\infty$-gömbbe való projekció alapvető helyességét igazolja.

\medskip
\noindent\rule{\linewidth}{0.3pt}

\medskip
\noindent\texttt{test\_pgd\_hits\_ball\_when\_step\_too\_large}: Szándékosan irreálisan nagy lépésméretet
($\alpha=10{,}0 \gg \epsilon=0{,}02$) ad meg, ami minden iterációban túllövést okoz; ilyenkor a projekciós
\texttt{clamp}-nek kell egyedül biztosítania a korlátot, ezért a tesztben a végső $|\delta|$ pontosan (nem
csak $\leq$) $\epsilon$-nak adódik minden pixelen -- ez a projekció, nem a lépésméret helyes megválasztása
általi korlátozottságot bizonyítja.

\medskip
\noindent\rule{\linewidth}{0.3pt}

\medskip
\noindent\texttt{test\_pgd\_clamps\_to\_valid\_pixel\_range}: A pixeltartomány-korlátozás ugyanúgy teljesül
PGD esetén is, mint FGSM-nél: szélsőséges (nulla) bemeneten és nagy $\epsilon$ mellett a kimenet a $[0, 1]$
tartományban marad.

\medskip
\noindent\rule{\linewidth}{0.3pt}

\medskip
\noindent\texttt{test\_pgd\_re\-sult\_has\_no\_graph\_history}: Ugyanaz az erős \texttt{detach}-bizonyíték
(\texttt{grad\_fn is None} és \texttt{is\_leaf is True}), mint FGSM esetén -- fontos, mert a PGD több
iterációt is végez, minden lépésben új \texttt{requires\_grad\_(True)} hívással, így könnyebben
,,elszivároghatna'' egy hivatkozás a gráfra, ha a ciklus végén elmaradna a leválasztás.

\medskip
\noindent\rule{\linewidth}{0.3pt}

\medskip
\noindent\texttt{test\_pgd\_ze\-ro\_epsi\-lon\_re\-turns\_original}: Határeset: $\epsilon=0$ esetén -- annak
ellenére, hogy több iteráció is lefut -- a projekció minden lépésben nulla szélességű intervallumra
szorítja a perturbációt, így a végeredmény megegyezik az eredeti képpel.

\medskip
\noindent\rule{\linewidth}{0.3pt}

\medskip
\noindent\texttt{test\_pgd\_un\-targeted\_re\-duces\_confidence\_in\_true\_label}: Viselkedési bizonyíték
\texttt{predict()}-en keresztül: a valódi címke konfidenciája több iteráció után is csökken -- ugyanaz az
elv, mint FGSM-nél, de itt azt is igazolja, hogy az iteratív frissítés nem ,,rontja el'' a kumulatív hatást.

\medskip
\noindent\rule{\linewidth}{0.3pt}

\medskip
\noindent\texttt{test\_pgd\_tar\-geted\_in\-creases\_confidence\_in\_target\_label}: A targeted PGD a célosztály
konfidenciáját növeli -- ugyanazon tulajdonságpár (non-targeted csökkenti/targeted növeli a releváns
konfidenciát) FGSM-re és PGD-re is parametrizált formában lett tesztelve
(\texttt{@pytest.mark.parametrize}), elkerülve, hogy csak az egyik algoritmusra íródjon le a viselkedés.

\bigskip

\section{Integrációs tesztek}

\subsection{GUI és támadási motor összekötése}

Az integrációs réteg feladata annak ellenőrzése, hogy az \texttt{app\_interface.py} eseménykezelői helyesen
építik-e fel az \texttt{AttackContext}-et a ténylegesen betöltött \texttt{ModelBundle}-ből, és a
\texttt{attack\_engine.attack()} hívás eredménye helyesen jelenik-e meg a felületen (kép- és
predikció-frissítés). Ez a réteg már nem mockolja a modellt, de a valódi felhasználói interakciót
(egérkattintás, csúszka-mozgatás) sem szimulálja még -- azt a GUI-tesztek rétege végzi.

\section{Automatizált GUI tesztek}

A GUI tesztek a felhasználói felületet magát, a ténylegesen megjelenő ablakot és widgeteket vizsgálják,
programozott módon szimulálva a felhasználói interakciókat, majd
ellenőrizi, hogy a felület a várt módon reagál-e. Ez a legfelső, egyben legköltségesebb és leglassabb
réteg a tesztpiramisban, hiszen minden egyes esethez fel kell építeni a teljes grafikus felületet -- ezért
csak kevés, a legfontosabb felhasználói forgatókönyveket lefedő tesztesetre érdemes szorítkozni, a
részletesebb logikai ellenőrzést az alsóbb rétegekre (egység- és integrációs tesztek) hagyva.

\section{Funkcionális tesztek és felhasználói validáció}

A grafikus felület funkcionális, ,,fekete doboz'' teszteseteit a felhasználói forgatókönyvek alapján,
\textit{Given--When--Then} formában fogalmazom meg. Az alábbi esetek a tényleges
felhasználói felületének elemeihez (gombok, csúszkák, célosztály-kereső) igazodnak, és kézi
(manuális) végrehajtásra tervezettek.

\begin{itemize}
  \item \textbf{Kép betöltése a ,,Load Image'' gombbal}\\
    AS A: felhasználó \textsc{I want to}: egy képet betölteni a vizsgálathoz\\
    \textsc{Given}: az alkalmazás elindítva, még nincs kép betöltve\\
    \textsc{When}: a ,,Load Image'' gombra kattint, majd egy érvényes képfájlt választ ki\\
    \textsc{Then}: az ,,Original'' képmezőben megjelenik a beolvasott, $224\times224$-re vágott
      kép, és alatta a top-3 predikció és konfidencia frissül

  \item \textbf{Modell váltása betöltött kép mellett}\\
    AS A: felhasználó \textsc{I want to}: összehasonlítani, hogy más architektúra hogyan látja a képet\\
    \textsc{Given}: egy kép már be van töltve, az ,,Original'' predikciós panel ki van töltve\\
    \textsc{When}: a modell-választó legördülőben ,,MobileNetV2''-t választ ,,ResNet50'' helyett\\
    \textsc{Then}: az ,,Original'' predikciók újraszámolódnak az új modellel; ha korábban már
      futott támadás, az ,,Adversarial'' predikciók is frissülnek

  \item \textbf{Non-targeted támadás végrehajtása}\\
    AS A: felhasználó \textsc{I want to}: általános tévesztést előidézni a modellben\\
    \textsc{Given}: kép betöltve, ,,Targeted attack'' jelölőnégyzet kikapcsolva, $\epsilon$ és
      iterációszám csúszkák beállítva\\
    \textsc{When}: az ,,Attack'' gombra kattint\\
    \textsc{Then}: a státuszsor ,,Running...'' majd ,,Attack complete'' üzenetet mutat, megjelenik
      a felnagyított perturbáció és a támadott kép, és az ,,Adversarial'' predikciós panelen az
      eredeti top-1 osztály konfidenciája jellemzően csökken

  \item \textbf{Targeted támadás célosztály kereséssel}\\
    AS A: felhasználó \textsc{I want to}: egy konkrét osztályra rávenni a modellt\\
    \textsc{Given}: kép betöltve, a célosztály-kereső mezőbe elkezd gépelni egy osztálynevet\\
    \textsc{When}: a szűrt listából kiválaszt egy osztályt, majd az ,,Attack'' gombra kattint\\
    \textsc{Then}: a ,,Targeted attack'' jelölőnégyzet automatikusan bejelölődik, a státuszsor a
      kiválasztott osztály nevét mutatja, és a támadás után a célosztály konfidenciája nő

  \item \textbf{Epsilon csúszka élő előnézete FGSM módban}\\
    AS A: felhasználó \textsc{I want to}: interaktívan látni, hogyan változik a perturbáció
      erőssége\\
    \textsc{Given}: FGSM módszerrel már lefutott egy támadás (a gradiens előjele cache-elve van)\\
    \textsc{When}: az $\epsilon$ csúszkát új értékre húzza, PGD-re váltás nélkül\\
    \textsc{Then}: a perturbáció és az ellenfél alapú kép új gradiens-számítás \emph{nélkül},
      azonnal újraszámolódik a cache-elt előjel alapján

  \item \textbf{Attack gomb kép nélkül}\\
    AS A: felhasználó \textsc{I want to}: (véletlenül) támadást indítani kép betöltése előtt\\
    \textsc{Given}: az alkalmazás elindítva, nincs betöltött kép\\
    \textsc{When}: az ,,Attack'' gombra kattint\\
    \textsc{Then}: egy ,,No image'' információs párbeszédablak jelenik meg, a támadás nem indul el

  \item \textbf{Védekezési szűrő alkalmazása (Blur / JPEG)}\\
    AS A: felhasználó \textsc{I want to}: megvizsgálni, hogy egy egyszerű szűrő semlegesíti-e a
      zajt\\
    \textsc{Given}: egy sikeres támadás után az ,,Adversarial'' kép elérhető\\
    \textsc{When}: a ,,Defend (Blur)'' vagy ,,Defend (JPEG)'' gombra kattint\\
    \textsc{Then}: az ,,Adversarial'' képmező a szűrt képre frissül, és a predikciós panel az
      újraosztályozott (a szűrés után kapott) eredményeket mutatja

  \item \textbf{Defend gomb ellenfél alapú kép nélkül}\\
    AS A: felhasználó \textsc{I want to}: (véletlenül) védekezést alkalmazni támadás előtt\\
    \textsc{Given}: kép be van töltve, de még nem futott támadás\\
    \textsc{When}: valamelyik ,,Defend'' gombra kattint\\
    \textsc{Then}: egy ,,No adversarial'' információs párbeszédablak jelenik meg, semmilyen kép
      nem változik

  \item \textbf{Reset gomb}\\
    AS A: felhasználó \textsc{I want to}: új kísérletet kezdeni ugyanazzal a képpel\\
    \textsc{Given}: egy támadás és esetleg egy védekezés már lefutott\\
    \textsc{When}: a ,,Reset'' gombra kattint\\
    \textsc{Then}: a perturbáció és az ,,Adversarial'' képmező üresre vált, az ,,Adversarial''
      predikciós panel névtelen (,,---'') állapotba áll vissza, az ,,Original'' kép és predikciói
      változatlanok maradnak

  \item \textbf{Hibakezelés támadás közben}\\
    AS A: felhasználó \textsc{I want to}: érthető visszajelzést kapni, ha a támadás valamiért
      meghiúsul\\
    \textsc{Given}: kép betöltve, a háttérszálon futó \texttt{\_attack\_worker} kivételt dob (pl.
      eszközhiba miatt)\\
    \textsc{When}: az ,,Attack'' gombra kattint\\
    \textsc{Then}: egy ,,Attack failed'' hibaablak jelenik meg a teljes veremkiíratással, a
      státuszsor ,,Attack failed.'' szöveget mutat, az alkalmazás nem fagy le
\end{itemize}

\section{Teszteredmények és értelmezésük}

\section{A tesztelés korlátai}

\chapter{Összefoglalás}

\section{Elért eredmények}

\section{Következtetések}

\section{Továbbfejlesztési lehetőségek}